\scanpage{82}
% archive foliation: 83

Si $k$ est un corps parfait, ceci donne une décomposition canonique en somme directe
\[
G\simeq G_{\mathrm{et}}\times G_{\mathrm{mult}}\times G_{\mathrm{bi}}
\]
et la catégorie des $G$ est can. \unsure{isomorphe} au produit des trois catégories correspondantes.

$G_{\mathrm{et}}$ et $G_{\mathrm{mult}}$, structure connue par théorie de Galois

$G_{\unsure{\mathrm{uni}|\mathrm{bi}}}$ ? Donné par théorie de Dieudonné.

\diagram{\begin{tikzcd}[column sep=small,row sep=small]
& \unsure{f_G} & |[draw,ellipse]|\text{nilpotent} \arrow[r,Leftrightarrow] \arrow[dll,no head,bend right=30] & G_{\mathrm{et}}=0\quad\text{i.e. }G^*\text{ unipotent} & f_G\text{ isom}\Longleftrightarrow G\text{ étale}\\
\parbox{3cm}{\margin{préciser que ça veut dire, en introduisant Frob. et Versch. \uline{itérés}.}} & w_G & |[draw,ellipse]|\text{nilpotent} \arrow[r,Leftrightarrow] \arrow[ll,no head,bend left=20] & G_{\mathrm{mult}}=0\quad\text{i.e. }G\text{ unipotent} & w_G\text{ isom}\Longleftrightarrow G\text{ d.m.}
\end{tikzcd}}

4. \uline{Théorie de Dieudonné}: \uline{le th. d'équivalence}.

On donne $k$ corps parfait de car. $p>0$

La théorie de Dieudonné fournit une \inserted{anti}équivalence
\diagram{\begin{tikzcd}[column sep=large]
p\text{-}\mathrm{Gr}_{\unsure{f}}(k) \arrow[r,"\approx"] &
\left\{\begin{array}{l}
\text{Catégorie des modules de longueur}\\
\text{finie }M\text{ sur }W,\text{ munis de}\\
F_M,V_M:M\mathbin{\unsure{-|\longrightarrow}}M\text{ satisfaisant}\\
F\lambda=\lambda^\sigma F,\quad\lambda V=V\lambda^\sigma\\
FV=p.\mathrm{id},\qquad FV=p.\mathrm{id}
\end{array}\right.\\
G \arrow[r,mapsto] & D^*(G)
\end{tikzcd}}

\uline{Interprétation}:
\diagram{\begin{tikzcd}
M^\sigma \arrow[r,shift left,"F_M"] & M \arrow[l,shift left,"V_M"]
\end{tikzcd}}

De plus, \unsure{\ldots}\ \deleted{\unsure{on a}}\deleted{\unsure{d'ici}} compte tenu que $D^*$ commute au changement de base, donc $D^*(G^{(p)})\simeq D^*(G)^{(p)}$ :
\[
F_{D^*(G)}=D^*(f_G),\qquad V_{D^*G}=D^{\unsure{*}}(w_G)
\]
d'où le dictionnaire :
\diagram{\left\{\begin{array}{lll}
G\text{ unipotent}&\Longleftrightarrow&V_M\text{ nilpotent}\\
G^*\text{ \unsure{unipotent}}&\Longleftrightarrow&F_M\text{ \unsure{unipotent|nilpotent}}\\
G\text{ étale}&\Longleftrightarrow&F_M\text{ isom}\\
G^*\text{ étale}&\Longleftrightarrow&V_M\text{ isom}\\
\text{i.e. }G,\text{ gr. d.m.}&&
\end{array}\right.}
